% foundations/metric_realization.tex

\chapter{Metric Realization}

\label{chap:metric-realization}

% ============================================================
\section{Motivation}
% ============================================================

\begin{remark}
Classical railway alignment systems often assume that geometry,
distance, curvature, projection, and realization are all defined within
one implicit Euclidean coordinate model.
\end{remark}

\begin{remark}
This assumption is sufficient for many engineering applications, but it
becomes problematic when:
\begin{itemize}
\item large-scale coordinate systems are involved,
\item Earth curvature becomes relevant,
\item projection distortions appear,
\item multiple CRS layers interact,
\item or alignment semantics must remain independent from metric
realization.
\end{itemize}
\end{remark}

\begin{remark}
Within AIM, a strict distinction is therefore introduced between:
\begin{itemize}
\item alignment semantics,
\item CRS embedding,
\item and metric realization.
\end{itemize}
\end{remark}

\begin{remark}
This separation enables:
\begin{itemize}
\item stable alignment semantics,
\item interchangeable metric models,
\item realization-aware geometry,
\item projection-aware interpretation,
\item and future differential-geometric extensions.
\end{itemize}
\end{remark}

% ============================================================
\section{Core Principle}
% ============================================================

\begin{principle}[Alignment semantics principle]
\label{princ:alignment-semantics}
The semantic structure of an alignment is independent of the metric
space in which it is realized.
\end{principle}

\begin{remark}
For example:
\begin{itemize}
\item straight elements,
\item constant-curvature arcs,
\item transition elements,
\item topology,
\item and switch logic
\end{itemize}
remain semantically stable even when the underlying metric model
changes.
\end{remark}

% ------------------------------------------------------------

\begin{principle}[Metric realization principle]
\label{princ:metric-realization}
Geometric realization depends on the metric model used for evaluation.
\end{principle}

\begin{remark}
Thus:
\begin{itemize}
\item distances,
\item curvature,
\item projection,
\item interpolation,
\item frame construction,
\item and geodesic behaviour
\end{itemize}
may differ between metric realizations.
\end{remark}

% ------------------------------------------------------------

\begin{principle}[Canonical separation principle]
\label{princ:canonical-metric-separation}
Alignment semantics and metric realization must remain conceptually
separable.
\end{principle}

% ============================================================
\section{Semantic Alignment Layer}
% ============================================================

\subsection{Intrinsic Semantics}

\begin{remark}
The semantic alignment layer describes:
\begin{itemize}
\item alignment topology,
\item curvature laws,
\item transition structure,
\item stationing semantics,
\item continuity conditions,
\item and sparse alignment structure.
\end{itemize}
\end{remark}

\begin{remark}
This layer is intentionally metric-independent.
\end{remark}

% ------------------------------------------------------------

\subsection{Sparse Representation}

\begin{remark}
Within AIM, the sparse alignment model stores:
\begin{itemize}
\item fixed curvature elements,
\item transition elements,
\item lookup-based curvature evolution,
\item sparse reconstruction parameters,
\item and associated semantic metadata.
\end{itemize}
\end{remark}

\begin{remark}
The sparse representation therefore acts as a semantic alignment
description rather than a finalized geometric embedding.
\end{remark}

% ============================================================
\section{CRS Embedding versus Metric Realization}
% ============================================================

\subsection{Conceptual Distinction}

\begin{remark}
Within AIM, CRS embedding and metric realization are related but
distinct concepts.
\end{remark}

\begin{remark}
The CRS layer defines:
\begin{itemize}
\item coordinate reference systems,
\item projection rules,
\item Earth-related embedding,
\item and coordinate transformations.
\end{itemize}
\end{remark}

\begin{remark}
The metric realization layer defines:
\begin{itemize}
\item distance interpretation,
\item curvature evaluation,
\item frame behaviour,
\item transport rules,
\item and geometric reconstruction behaviour.
\end{itemize}
\end{remark}

% ------------------------------------------------------------

\subsection{Coordinate Transformation}

\begin{remark}
Coordinate transformations are represented by:
\[
\mathcal C
:
\Omega_{\mathrm{world}}
\rightarrow
\Omega_{\mathrm{world}}.
\]
\end{remark}

\begin{remark}
Coordinate transformation alone does not define intrinsic railway
geometry.
\end{remark}

% ------------------------------------------------------------

\subsection{Projection Dependence}

\begin{proposition}
\label{prop:projection-distorts-alignment}
A straight line in one projected CRS is generally not identical to a
straight line in another projection space.
\end{proposition}

\begin{remark}
Similarly:
\begin{itemize}
\item curvature,
\item transition behaviour,
\item stationing interpretation,
\item and lateral offsets
\end{itemize}
may become projection-dependent.
\end{remark}

% ============================================================
\section{Metric Realization Spaces}
% ============================================================

\subsection{General Interpretation}

\begin{definition}[Metric realization space]
\label{def:metric-space}

A metric realization space defines:
\begin{itemize}
\item distance interpretation,
\item directional interpretation,
\item curvature evaluation,
\item frame transport,
\item projection behaviour,
\item and geometric reconstruction rules.
\end{itemize}

\end{definition}

% ------------------------------------------------------------

\subsection{Metric Realization Operator}

\begin{definition}[Metric realization operator]
\label{def:metric-realization-operator}

Metric realization is represented abstractly by:
\[
\mathcal G
:
\mathcal X
\rightarrow
\mathcal X_{\mathrm{metric}}.
\]

\end{definition}

% ------------------------------------------------------------

\begin{remark}
The operator \(\mathcal G\) defines how semantic railway states are
evaluated geometrically within a selected metric model.
\end{remark}

% ------------------------------------------------------------

\subsection{Operator Boundary}

\begin{principle}[Metric operator boundary]
\label{princ:metric-operator-boundary}
Alignment logic should not directly assume a specific metric model.
\end{principle}

\begin{remark}
Instead, alignment operations interact with geometry through metric
operators and runtime evaluation services.
\end{remark}

\begin{remark}
This creates a strict separation between:
\begin{itemize}
\item semantic alignment structure,
\item and metric realization behaviour.
\end{itemize}
\end{remark}

% ============================================================
\section{Engineering Metric}
% ============================================================

\subsection{Role}

\begin{remark}
The dominant practical realization mode within AIM is the engineering
metric.
\end{remark}

% ------------------------------------------------------------

\subsection{Characteristics}

\begin{remark}
The engineering metric assumes:
\begin{itemize}
\item locally Euclidean geometry,
\item Cartesian coordinates,
\item fast evaluation,
\item simple frame construction,
\item and numerically stable reconstruction.
\end{itemize}
\end{remark}

% ------------------------------------------------------------

\subsection{Importance}

\begin{principle}[Engineering priority principle]
\label{princ:engineering-priority}
The engineering metric remains the primary operational mode of AIM and
must not suffer unnecessary computational overhead.
\end{principle}

\begin{remark}
Advanced metric concepts must therefore remain compatible with highly
efficient engineering-scale evaluation.
\end{remark}

% ============================================================
\section{Projection-Aware Metric}
% ============================================================

\subsection{Motivation}

\begin{remark}
Projected coordinate systems introduce:
\begin{itemize}
\item scale distortions,
\item directional distortions,
\item projection-dependent curvature effects,
\item and coordinate-dependent metric behaviour.
\end{itemize}
\end{remark}

% ------------------------------------------------------------

\subsection{Interpretation}

\begin{remark}
Within AIM, projection-aware realization interprets alignment geometry
through projection-dependent metric evaluation.
\end{remark}

\begin{remark}
Thus, runtime geometry depends not only on intrinsic alignment semantics
but also on projection context.
\end{remark}

% ------------------------------------------------------------

\subsection{Operational Consequences}

\begin{remark}
Projection-aware realization affects:
\begin{itemize}
\item world--track projection,
\item nearest-point search,
\item stationing interpretation,
\item local frame construction,
\item and geometric reconstruction.
\end{itemize}
\end{remark}

% ============================================================
\section{Ellipsoid-Based Metric}
% ============================================================

\subsection{Motivation}

\begin{remark}
At sufficiently large spatial scales, Earth curvature becomes relevant
for railway alignment interpretation.
\end{remark}

% ------------------------------------------------------------

\subsection{Geodesic Interpretation}

\begin{remark}
Within ellipsoid-based geometry:
\begin{itemize}
\item Euclidean straightness,
\item geodesic straightness,
\item and engineering straightness
\end{itemize}
are no longer identical concepts.
\end{remark}

% ------------------------------------------------------------

\subsection{Curvature Interpretation}

\begin{remark}
On curved surfaces, curvature decomposes into:
\[
\kappa^2
=
\kappa_g^2
+
\kappa_n^2,
\]
where:
\begin{itemize}
\item \(\kappa_g\) denotes geodesic curvature,
\item \(\kappa_n\) denotes surface-normal curvature.
\end{itemize}
\end{remark}

\begin{remark}
For railway dynamics, geodesic curvature is generally the physically
relevant quantity.
\end{remark}

% ------------------------------------------------------------

\subsection{Frame Interpretation}

\begin{remark}
Within ellipsoid-aware geometry, frame construction may require:
\begin{itemize}
\item tangent transport,
\item surface-normal evaluation,
\item and connection-dependent orientation rules.
\end{itemize}
\end{remark}

% ============================================================
\section{Future Differential-Geometric Layer}
% ============================================================

\subsection{Motivation}

\begin{remark}
Future railway alignment systems may require:
\begin{itemize}
\item intrinsic surface geometry,
\item Earth-attached frames,
\item geodesic transition interpretation,
\item and fully differential-geometric reconstruction methods.
\end{itemize}
\end{remark}

% ------------------------------------------------------------

\subsection{Operator-Based Geometry}

\begin{remark}
Within such a framework, geometric operations would no longer be based
primarily on Cartesian coordinate algebra.
\end{remark}

\begin{remark}
Instead, geometry would emerge from:
\begin{itemize}
\item local operators,
\item metric tensors,
\item connection rules,
\item covariant derivatives,
\item and differential-geometric transport.
\end{itemize}
\end{remark}

% ============================================================
\section{Runtime Dependency}
% ============================================================

\subsection{Runtime Evaluation}

\begin{remark}
Metric realization affects runtime services such as:
\begin{itemize}
\item projection,
\item stationing,
\item frame construction,
\item nearest-point search,
\item pose3 evaluation,
\item and geometric reconstruction.
\end{itemize}
\end{remark}

% ------------------------------------------------------------

\subsection{Metric-Aware Runtime}

\begin{remark}
Runtime systems should therefore evaluate geometry through metric-aware
operator services rather than through hardcoded Euclidean assumptions.
\end{remark}

% ------------------------------------------------------------

\subsection{pose3 Dependency}

\begin{remark}
Runtime pose3 evaluation depends jointly on:
\begin{itemize}
\item semantic alignment states,
\item metric realization,
\item CRS embedding,
\item frame operators,
\item profile interpretation,
\item and cant interpretation.
\end{itemize}
\end{remark}

\begin{remark}
Thus, pose3 is not a purely intrinsic object, but a runtime-generated
metric interpretation of semantic alignment structure.
\end{remark}

% ============================================================
\section{Conceptual Architecture}
% ============================================================

\begin{figure}[h]
\centering

\begin{tikzpicture}[
node distance=1.5cm and 2.2cm,
box/.style={
draw,
rounded corners,
align=center,
minimum width=3.7cm,
minimum height=1.0cm
},
arrow/.style={->, thick}
]

\node[box] (semantic) {
Alignment Semantics\\
Sparse Alignment\\
Transitions / Stationing
};

\node[box, below=of semantic] (metric) {
Metric Realization Layer\\
Operator Boundary\\
$\mathcal G$
};

\node[box, below left=of metric] (euclid) {
Engineering Metric\\
Euclidean / Fast
};

\node[box, below=of metric] (proj) {
Projection-Aware Metric
};

\node[box, below right=of metric] (ellipsoid) {
Ellipsoid Metric
};

\node[box, below=2.2cm of proj] (runtime) {
Runtime Evaluation\\
pose3 / Projection / Frames
};

\draw[arrow] (semantic) -- (metric);

\draw[arrow] (metric) -- (euclid);
\draw[arrow] (metric) -- (proj);
\draw[arrow] (metric) -- (ellipsoid);

\draw[arrow] (euclid) -- (runtime);
\draw[arrow] (proj) -- (runtime);
\draw[arrow] (ellipsoid) -- (runtime);

\end{tikzpicture}

\caption{
AIM separates semantic alignment structure from metric realization.
Different metric backends may realize the same semantic alignment model,
while runtime services operate on the resulting metric interpretation.
}
\label{fig:metric-realization-architecture}

\end{figure}

% ============================================================
\section{Key Insight}
% ============================================================

\begin{proposition}
Railway alignment semantics and metric realization are fundamentally
different conceptual layers.
\end{proposition}

\begin{proposition}
The same semantic alignment structure may be realized within different
metric spaces.
\end{proposition}

\begin{proposition}
Metric realization therefore acts as an operator layer between semantic
alignment logic and spatial geometry evaluation.
\end{proposition}

\begin{proposition}
Runtime geometry emerges from interaction between:
\begin{itemize}
\item semantic alignment structure,
\item metric realization,
\item CRS embedding,
\item and runtime evaluation operators.
\end{itemize}
\end{proposition}

% ============================================================
\section{Connection to AIM}
% ============================================================

\begin{summary}
Within AIM, alignment semantics are intentionally separated from metric
realization.

This enables:
\begin{itemize}
\item stable sparse alignment representations,
\item interchangeable metric models,
\item realization-aware geometry,
\item projection-aware evaluation,
\item runtime-dependent pose3 interpretation,
\item and future differential-geometric extensions.
\end{itemize}

The metric realization layer therefore forms the operator boundary
between semantic alignment logic and operational geometry.
\end{summary}
